% ============================================================================
%  VSEPR-SIM — §P: Petalized Compute Allocation Model
%  Resource Partitioning as Botanical Abstraction.
%  Printer-ready LaTeX.  Compile with pdflatex (twice for refs).
% ============================================================================
\documentclass[11pt, a4paper, oneside]{article}

% ── Geometry ────────────────────────────────────────────────────────────────
\usepackage[
  top=1.0in, bottom=1.0in, left=1.15in, right=1.15in,
  headheight=14pt
]{geometry}

% ── Fonts / encoding ───────────────────────────────────────────────────────
\usepackage[T1]{fontenc}
\usepackage{lmodern}
\usepackage{microtype}

% ── Math ───────────────────────────────────────────────────────────────────
\usepackage{amsmath, amssymb, amsthm, mathtools}
\usepackage{bm}

% ── Tables ─────────────────────────────────────────────────────────────────
\usepackage{booktabs}
\usepackage{array}
\usepackage{tabularx}
\newcolumntype{L}[1]{>{\raggedright\arraybackslash}p{#1}}
\newcolumntype{C}[1]{>{\centering\arraybackslash}p{#1}}

% ── Lists ──────────────────────────────────────────────────────────────────
\usepackage{enumitem}
\setlist{nosep, leftmargin=1.5em}

% ── Code listings ─────────────────────────────────────────────────────────
\usepackage{listings}
\usepackage{xcolor}
\definecolor{lstbg}{RGB}{248,248,248}
\definecolor{lstkw}{RGB}{0,0,180}
\definecolor{lstcm}{RGB}{80,120,80}
\lstdefinestyle{cppstyle}{
  language=C++, basicstyle=\ttfamily\footnotesize,
  backgroundcolor=\color{lstbg},
  keywordstyle=\color{lstkw}\bfseries,
  commentstyle=\color{lstcm}\itshape,
  frame=single, framesep=3pt, breaklines=true, keepspaces=true, tabsize=4,
  showstringspaces=false
}
\lstset{style=cppstyle}

% ── Theorem environments ───────────────────────────────────────────────────
\theoremstyle{definition}
\newtheorem{definition}{Definition}[section]
\newtheorem{postulate}[definition]{Postulate}
\newtheorem{remark}[definition]{Remark}

% ── TColorbox ─────────────────────────────────────────────────────────────
\usepackage{tcolorbox}
\tcbuselibrary{skins, breakable}

\definecolor{rootcol}{RGB}{235,225,200}
\definecolor{stemcol}{RGB}{230,245,220}
\definecolor{corecol}{RGB}{255,245,220}
\definecolor{petalcol}{RGB}{255,230,240}
\definecolor{bloomcol}{RGB}{230,240,255}

% ── Headers / footers ─────────────────────────────────────────────────────
\usepackage{fancyhdr}
\pagestyle{fancy}
\fancyhf{}
\fancyhead[L]{\small\textit{VSEPR-SIM Research Platform}}
\fancyhead[R]{\small\textit{§P — Petalized Compute Allocation Model}}
\fancyfoot[C]{\thepage}
\renewcommand{\headrulewidth}{0.4pt}

% ── Section formatting ─────────────────────────────────────────────────────
\usepackage{titlesec}
\titleformat{\section}{\Large\bfseries}{§\thesection}{1em}{}
\titleformat{\subsection}{\large\bfseries}{\thesubsection}{1em}{}
\titleformat{\subsubsection}{\normalsize\bfseries}{\thesubsubsection}{1em}{}

% ── Hyperlinks ─────────────────────────────────────────────────────────────
\usepackage{hyperref}
\hypersetup{colorlinks=true,linkcolor=blue!60!black,urlcolor=blue!70!black}

% ── Macros ─────────────────────────────────────────────────────────────────
\newcommand{\file}[1]{\texttt{#1}}
\newcommand{\type}[1]{\texttt{#1}}

% ============================================================================
\begin{document}
% ============================================================================

\begin{titlepage}
\centering
\vspace*{3cm}

{\Huge\bfseries Petalized Compute\\[6pt] Allocation Model\par}
\vspace{1.0cm}
{\Large\itshape Resource Partitioning as Botanical Abstraction\par}
\vspace{1.5cm}
{\large VSEPR-SIM Research Platform\par}
\vspace{0.5cm}
{\large Section §P — Theoretical Foundation\par}
\vspace{2.0cm}
{\normalsize Version 3.0.0 — First Principles Draft\par}
\vfill
{\small\today}
\end{titlepage}

\tableofcontents
\newpage

% ============================================================================
\section{Motivation: The Botanical Metaphor for Computation}
% ============================================================================

Modern compute environments expose resources as anonymous pools: a count of
cores, a quantity of memory, a fractional share of accelerator time.  While
technically sufficient, this flat representation obscures the \emph{structure}
of resource availability and the \emph{dependency relationships} that connect
visible computation to its supporting infrastructure.

The \textbf{Petalized Compute Allocation Model} introduces a botanical
metaphor that maps resource partitioning onto the anatomy of a flower.
The metaphor is not decorative; it encodes a specific hierarchy:

\begin{itemize}
  \item The \textbf{bloom} is the exposed computational surface of the machine.
  \item Each \textbf{petal} is a schedulable compute partition.
  \item The \textbf{roots} and \textbf{stem} represent the non-visible support
        systems required to sustain execution: power delivery, thermal
        stability, storage movement, network access, and system coordination.
\end{itemize}

The metaphor thus captures both the \emph{beauty of variation} (heterogeneous
resource profiles across partitions) and the \emph{physical dependency} of
computation on deeper infrastructural support.

% ============================================================================
\section{Formal Flower Decomposition}
% ============================================================================

\begin{definition}[Flower]
A \textbf{flower} $F$ is a quadruple
\begin{equation}
  F = \bigl\{ R,\; S,\; C,\; P \bigr\}
\end{equation}
where each component represents a distinct functional layer of the compute
system.
\end{definition}

\begin{tcolorbox}[colback=rootcol, title={\textbf{$R$ — Root Systems (Support Layer)}}]
Power delivery, cooling, storage, network ingress, persistence.
The root systems are the non-visible infrastructure upon which all
computation ultimately depends.  A root failure collapses the entire flower.
\end{tcolorbox}

\begin{tcolorbox}[colback=stemcol, title={\textbf{$S$ — Stem (Transport and Control Backbone)}}]
Scheduler path, process routing, kernel coordination.
The stem connects the exposed bloom to the root systems, providing the
transport and arbitration layer through which resource requests propagate.
\end{tcolorbox}

\begin{tcolorbox}[colback=corecol, title={\textbf{$C$ — Core (Orchestration Center)}}]
Allocator, supervisor, node manager, system controller.
The core is the decision point: it determines how resources are distributed
across petals and how competing demands are resolved.
\end{tcolorbox}

\begin{tcolorbox}[colback=petalcol, title={\textbf{$P$ — Petals (Exposed Compute Partitions)}}]
The set of schedulable compute partitions:
\begin{equation}
  P = \bigl\{ p_1,\; p_2,\; \ldots,\; p_n \bigr\}
\end{equation}
Each petal is a bounded, explicitly described region of machine capability
visible to user-facing workloads.
\end{tcolorbox}

% ============================================================================
\section{Petal Structure}
% ============================================================================

\begin{definition}[Petal]
Each petal $p_i$ is defined as
\begin{equation}
  p_i = \bigl\{ M_i,\; C_i,\; G_i,\; D_i \bigr\}
\end{equation}
where:
\begin{itemize}
  \item $M_i$ — memory assigned to petal $i$ (e.g.\ in GB),
  \item $C_i$ — CPU cores assigned to petal $i$,
  \item $G_i$ — GPU fraction or accelerator share assigned to petal $i$
        ($0 \le G_i \le 1$),
  \item $D_i$ — number of internal divisions (sub-partitions) available
        within petal $i$.
\end{itemize}
\end{definition}

\begin{remark}[Sub-division]
The internal divisions $D_i$ allow tasks to be distributed at finer
granularity within the same visible compute region.  A petal with $D_i = 4$
can host up to four independently scheduled work units, each receiving an
equal share of $M_i$, $C_i$, and $G_i$ unless explicitly overridden by the
orchestration core.
\end{remark}

\subsection{Concrete Example}

Consider a single node represented as a flower whose petals each expose:

\begin{equation}
  p_i = \bigl\{ 25.6\;\text{GB},\; 6\;\text{cores},\; 0.30\;\text{GPU},\; 4 \bigr\}
\end{equation}

This states that petal $i$ provides 25.6\,GB of memory, 6~CPU cores, 30\% of
a GPU allocation, and 4~sub-divisions for internal scheduling.

% ============================================================================
\section{Bloom Capacity}
% ============================================================================

\begin{definition}[Bloom Capacity]
The \textbf{total exposed bloom capacity} $B$ of a flower with $n$ petals is
\begin{equation}
  B = \sum_{i=1}^{n} p_i
\end{equation}
interpreted component-wise:
\begin{equation}
  B = \left(
    \sum_{i=1}^{n} M_i,\;\;
    \sum_{i=1}^{n} C_i,\;\;
    \sum_{i=1}^{n} G_i,\;\;
    \sum_{i=1}^{n} D_i
  \right)
\end{equation}
\end{definition}

For a flower with $n = 4$ identical petals as in the example above:

\begin{equation}
  B = \bigl( 102.4\;\text{GB},\; 24\;\text{cores},\; 1.20\;\text{GPU},\; 16 \bigr)
\end{equation}

Note that $\sum G_i > 1$ is valid when a node contains multiple physical
accelerators or when GPU fractions are expressed relative to a single device
among several.

% ============================================================================
\section{Heterogeneous Petal Configurations}
% ============================================================================

The petalized model is particularly suitable for heterogeneous or mixed-load
systems.  Petals need not be uniform:

\begin{center}
\begin{tabular}{@{} l C{2.2cm} C{1.8cm} C{1.8cm} C{1.4cm} @{}}
\toprule
\textbf{Petal} & \textbf{Memory (GB)} & \textbf{CPU Cores} & \textbf{GPU Share} & \textbf{Divisions} \\
\midrule
$p_1$ (CPU-dense)     & 16.0  & 12 & 0.00 & 6 \\
$p_2$ (Memory-heavy)  & 64.0  & 4  & 0.00 & 2 \\
$p_3$ (Accelerated)   & 8.0   & 2  & 0.50 & 1 \\
$p_4$ (Balanced)      & 25.6  & 6  & 0.30 & 4 \\
\bottomrule
\end{tabular}
\end{center}

Color, placement, and visual encoding may then represent:
\begin{itemize}
  \item node role (compute, storage, accelerator),
  \item priority class (interactive, batch, preemptible),
  \item thermal state (nominal, throttled, critical),
  \item resource pressure (idle, loaded, saturated).
\end{itemize}

A petal is therefore both a \textbf{compute unit} and a \textbf{visual
contract}: it tells the user what resources exist, how they are divided, and
where tasks are currently placed.

% ============================================================================
\section{Scheduling Interpretation}
% ============================================================================

The flower decomposition provides a natural mapping to scheduling abstractions:

\begin{enumerate}
  \item \textbf{Petal selection} — the scheduler selects a petal $p_i$ whose
        resource profile matches the task requirements.
  \item \textbf{Division assignment} — within the selected petal, the task
        is assigned to one of $D_i$ available sub-partitions.
  \item \textbf{Resource binding} — the task receives its share of $M_i$,
        $C_i$, and $G_i$ proportional to the division count.
  \item \textbf{Bloom monitoring} — the orchestration core $C$ monitors
        aggregate bloom utilization and may rebalance across petals.
\end{enumerate}

This hierarchy maps cleanly onto existing container orchestration, NUMA-aware
scheduling, and GPU partitioning (e.g.\ MIG, MPS) paradigms while providing a
human-readable structural overlay.

% ============================================================================
\section{Infrastructure Dependencies}
% ============================================================================

The root and stem layers encode a critical design principle: \emph{visible
computation depends on invisible infrastructure}.  A flower whose roots are
compromised (e.g.\ insufficient cooling, failing storage) cannot sustain its
bloom regardless of how well petals are configured.

\subsection{Root System Components}

\begin{itemize}
  \item \textbf{Power delivery} — wattage budget, PSU redundancy, UPS state.
  \item \textbf{Cooling} — ambient temperature, fan RPM, thermal throttle
        thresholds.
  \item \textbf{Storage} — IOPS capacity, read/write bandwidth, persistence
        guarantees.
  \item \textbf{Network ingress} — link bandwidth, latency, packet loss rate.
\end{itemize}

\subsection{Stem Components}

\begin{itemize}
  \item \textbf{Scheduler path} — queue depth, dispatch latency, fairness
        policy.
  \item \textbf{Process routing} — affinity masks, NUMA topology, interrupt
        steering.
  \item \textbf{Kernel coordination} — syscall overhead, context-switch cost,
        lock contention.
\end{itemize}

% ============================================================================
\section{Thermal Monitoring Integration}
% ============================================================================

The root layer naturally incorporates thermal monitoring as a first-class
concern.  Temperature readings from hardware sensors can be converted to
a consistent Celsius representation:

\begin{equation}
  T_{\text{C}} = \frac{T_{\text{raw}}}{10} - 273.15
\end{equation}

where $T_{\text{raw}}$ is the value reported by firmware interfaces such as
\type{MsAcpi\_ThermalZoneTemperature} on Windows systems.  Thermal state
feeds back into petal scheduling: a thermally stressed root may force the
orchestration core to reduce bloom capacity by disabling or throttling
individual petals.

% ============================================================================
\section{GPU Telemetry as Petal Metadata}
% ============================================================================

For accelerator-bearing petals, runtime telemetry enriches the petal
description with dynamic state:

\begin{itemize}
  \item \textbf{GPU utilization} $u_{\text{gpu}}$ — fraction of active SMs
        (0--100\%).
  \item \textbf{VRAM usage} $m_{\text{gpu}}$ — memory currently allocated
        (in MB or GB).
\end{itemize}

These values are cached to avoid per-prompt overhead (e.g.\ querying
\type{nvidia-smi} on every render cycle).  A cache lifetime of
$\sim$\!2\,s balances freshness against launch cost.

Color encoding of GPU state within the petal visualization follows the
convention:

\begin{center}
\begin{tabular}{@{} l l l @{}}
\toprule
\textbf{Metric} & \textbf{Threshold} & \textbf{Color} \\
\midrule
$u_{\text{gpu}} \ge 80\%$  & High load    & Red (ANSI 196) \\
$u_{\text{gpu}} \ge 50\%$  & Medium load  & Orange (ANSI 208) \\
$u_{\text{gpu}} < 50\%$    & Low load     & Green (ANSI 118) \\
$m_{\text{gpu}}$            & Any          & Light cyan (ANSI 159) \\
\bottomrule
\end{tabular}
\end{center}

% ============================================================================
\section{Relationship to the VSEPR-SIM Platform}
% ============================================================================

Within VSEPR-SIM, the petalized compute allocation model serves as the
conceptual bridge between the atomistic simulation kernel and the physical
hardware that executes it.  The mapping is:

\begin{itemize}
  \item \textbf{Atomistic structure generation} tasks are assigned to petals
        based on memory and core requirements.
  \item \textbf{Visualization rendering} tasks may require accelerator-bearing
        petals with nonzero $G_i$.
  \item \textbf{Reporting and export} tasks are lightweight and may occupy
        minimal petals or share divisions within a larger petal.
  \item \textbf{Discovery workflows} involving parameter sweeps distribute
        candidate evaluations across the full bloom $B$.
\end{itemize}

This ensures that resource allocation is explicit, inspectable, and
deterministic — consistent with the anti-black-box philosophy of the
platform.

% ============================================================================
\section{Summary}
% ============================================================================

The Petalized Compute Allocation Model provides:

\begin{enumerate}
  \item A structured decomposition of compute resources into a four-layer
        botanical hierarchy: roots, stem, core, and petals.
  \item A formal definition of each petal as an explicit tuple of memory,
        CPU, GPU, and sub-division capacity.
  \item A component-wise bloom capacity metric that aggregates across all
        petals.
  \item Support for heterogeneous configurations where petals vary in
        resource profile and scheduling role.
  \item Natural integration with thermal monitoring, GPU telemetry, and
        infrastructure health as feedback signals.
  \item A human-readable visual contract that makes resource state
        inspectable at every level.
\end{enumerate}

The model is both a \textbf{scheduling abstraction} and an
\textbf{interface abstraction}: it organizes computation into explicit,
auditable partitions while providing a visual language for communicating
system state to the researcher.

% ============================================================================
\end{document}
% ============================================================================
